% Nota de onde_marcar. GERADA por flim_al/tabelas.py.
\footnotesize Orçamento em \textbf{pixels anotados por imagem}, não em imagens: o que limita o FLIM é a quantidade de pixels marcados --- a mesma imagem rende 54 kernels por camada com traço esparso e 200 com traço denso. Todos os braços gastam exatamente o mesmo número de pixels, nas MESMAS imagens (as do artigo); o que muda é onde eles caem. \texttt{uniforme} é o FLIM de hoje. \texttt{uniforme\_balanceado} tem a mesma proporção objeto/fundo que a incerteza produz, com os pixels sorteados --- é ele que separa *onde se anotou* de *quanto de cada classe se anotou*, porque escolher região incerta escolhe junto muito mais foreground (~73\% contra ~16\%, medido antes de rodar). \texttt{regiao\_borda} usa o gabarito e é teto, não método. Diferenças pareadas por semente; p de t pareado bicaudal. Campanha \texttt{onde\_marcar/k3/teste60/seg300}.
\normalsize
